\section{认知智能：感知智能之后的下一步}

2015--2016 年，上一代 AI 公司主要围绕 CV、早期 NLP 和深度学习落地。张鹏说，当时他们意识到感知智能有天花板：人脸识别超过人以后，再往上提升意义有限，产业效果也接近上限。于是他们提出“认知智能”，作为迈向通用人工智能的下一层台阶。

\begin{figure}[H]
\centering
\includegraphics[width=0.94\textwidth]{cognitive-intelligence-ladder.png}
\caption{从感知智能到认知智能，再到大模型、Agent 和 AGI。}
\end{figure}

\begin{knowledgebox}{读图：认知智能是阶段性定义}
图中认知智能不是 AGI 本身，而是感知智能之后的 next step。它试图把 AI 从识别、分类、感知任务推进到理解、推理、知识、语言和复杂任务。后来大模型成为承载这个方向的重要技术形态。
\end{knowledgebox}

\subsection{为什么不直接定义 AGI}

张鹏说，当时他们并不强行定义 AGI，因为太遥远、也定义不清楚；但“下一步是什么”可以尝试定义。这个做法很稳健：不要用终极概念替代阶段性问题。认知智能为智谱提供了一个足够远、但仍能工程化推进的方向。

\begin{importantbox}{阶段性目标的价值}
长期愿景需要阶段性定义。AGI 是远目标，认知智能是可讨论、可组织资源、可做工程转化的阶段目标。
\end{importantbox}

\subsection{本章小结}

认知智能是智谱早期的方向锚点。它让公司既不止步于感知智能，也不空喊 AGI，而是寻找从理解、推理和知识工程进入下一代 AI 的路线。

